\chapter{Introduction}

Finite-state devices are fundamental to computer science, finding applications in
compiler design \cite{aho_lam_sethi_ullman}, natural language processing
\cite{mohri_applications}, and other fields. An introduction to finite-state
automata and their relation to the theory of computation is given in
\cite{hopcroft_motwani_ullman}. A generalized model of finite-state automata,
finite-state transducers, and bimachines is presented in \cite{mihov_schulz}.
\bigbreak
A central problem in the theory of finite-state devices is that of non-determinism
and the existence of deterministic representations with similar expressive power to
that of non-deterministic machines.
In their seminal work \cite{rabin_scott}, Rabin and Scott introduce the notion of
non-deterministic finite-state automata, show that these devices recognize precisely
the languages recognized by deterministic finite-state automata, and present an
algorithm that, given a non-deterministic finite-state automaton, constructs a
deterministic finite-state automaton recognizing the same language. Mohri shows that
this determinization algorithm can be generalized to rational functions of bounded
variation \cite{mohri_applications}. Beal et al. present simple algorithms for
deciding whether the relation represented by a finite-state transducer is a function
of bounded variation \cite{beal}. While often good in practice, a downside to the
determinization algorithm is that in the worst case it produces a deterministic
finite-state device with an exponential number of states compared to the input
non-deterministic machine.
\bigbreak
A concept relevant to the representation of a finite-state device is that of a graph.
Algorithms over graphs are studied outside the theory of finite-state devices but in
many cases find applications to related problems. A class of efficient algorithms for
graphs that contain a separator is based on the divide-and-conquer method
\cite{aho_hopcroft}. Lipton and Tarjan show that every planar graph contains a
separator of relatively small size \cite{lipton_tarjan_theorem}, making these
algorithms applicable to many practical problems \cite{lipton_tarjan_applications}.
Similar results follow for other classes of graphs \cite{gilbert_hutchinson_tarjan}
\cite{alon_seymour_thomas}.
\bigbreak
In this thesis, we take a closer look at graphs corresponding to finite-state
automata and finite-state transducers, and show how a separator in the graph can be
used to obtain deterministic representations that are smaller than those produced by
the classical constructions. If a set of states $X$ splits the remaining states into
components $Q_1, \ldots, Q_k$ that are not connected by transitions, a
\emph{quasi-deterministic} finite-state automaton\footnote{The term is also used by
Caron, Flouret and Mignot \cite{caron_flouret_mignot} for a different notion:
structures derived from $(k,l)$-unambiguous automata that use a window of fixed size.
The two notions are unrelated.} keeps one deterministic quasi-state per component and
synchronizes the components through the separator. The number of components controls a
trade-off. With a single component, a quasi-deterministic finite-state automaton is a
deterministic finite-state automaton: it performs one look-up per input symbol but may
need up to $2^{|Q|}$ states. With $k$ components, it performs $\mathcal{O}(k)$ look-ups
per input symbol and has at most $2^{|Q_i| + |X|}$ quasi-states in its $i$-th
component.
\bigbreak
The main contributions of this thesis are the following.
\begin{itemize}
  \item A quasi-determinization construction for $k$-separated finite-state automata,
    together with a characterization of its reachable states: they are exactly the
    projections of the states of the subset construction onto the components
    (\Cref{lem:quasidet_invariant}).
  \item An analysis of the construction time. A direct construction processes as many
    states as the subset construction (\Cref{prop:quasidet_iterations}). We give a
    per-component construction whose running time depends only on the sizes of the
    components and the separator (\Cref{thm:quasidet_components}). Both the size and
    the running time depend only on the essential states of the components: the start
    states and the targets of non-$\epsilon$ transitions
    (\Cref{lem:quasidet_essential,thm:quasidet_components_essential}).
  \item Separators of size $2$ for finite-state automata constructed from regular
    expressions, computed in linear time, whose components have at most
    $\frac{2}{3}|Q|$ states (\Cref{thm:thompson_separator}) or at most $\frac{2}{3}m +
    1$ essential states, where $m$ is the number of symbol occurrences
    (\Cref{thm:thompson_separator_symbols}). As a consequence, every regular expression
    can be quasi-determinized in $\mathcal{O}(|\Sigma| |Q| 2^{\frac{2}{3}m})$ time
    (\Cref{cor:quasidet_regex}), whereas the subset construction may produce up to
    $2^{m+1}$ states.
  \item Quasi-subsequential finite-state transducers and a construction of such a
    transducer for every trimmed functional $k$-separated finite-state transducer
    whose function has bounded variation
    (\Cref{chapter:separators_finite_state_transducers}).
  \item A formalization of the main results in the Lean theorem prover, and an
    implementation of the constructions for finite-state automata with an
    experimental comparison.
\end{itemize}
\Cref{chapter:preliminaries} introduces the notation.
\Cref{chapter:separators_finite_state_automata} studies separated finite-state
automata, \Cref{chapter:separators_finite_state_transducers} separated finite-state
transducers, and \Cref{chapter:conclusion} summarizes the results and discusses future
work.

\section{Related work}

\paragraph{Decompositions of automata constructed from regular expressions.}
The observation that the automata constructed from regular expressions by the algorithm
of Thompson \cite{thompson} decompose along the parse tree of the expression into parts
that share only their start and final states is not new. Myers \cite{myers} uses it to
match regular expressions faster than by simulating the nondeterministic automaton: the
automaton is decomposed into parts of logarithmic size, and for each part transitions
between sets of its states are tabulated in advance. Bille and Farach-Colton
\cite{bille_farach_colton} refine this approach and reduce its space requirements. In
these algorithms, the parts are small, their number grows with the size of the
expression, and the tables are used to simulate the nondeterministic automaton on a
text. In contrast, we consider arbitrary separated finite-state automata, construct a
quasi-deterministic finite-state automaton in advance whose reachable states correspond
exactly to the states of the subset construction, and study how its size and the cost
of its construction depend on the separator and the components. With one component,
the construction yields a deterministic finite-state automaton; decompositions into
many small parts are closer in spirit to the tabulation approach of Myers.

\paragraph{Partial determinization.}
In practice, the cost of determinization is often reduced by not determinizing the whole
automaton at once. Van Noord \cite{van_noord} compares variants of the subset
construction for automata with many $\epsilon$-transitions, as they arise in natural
language processing. Becchi and Crowley \cite{becchi_crowley} keep the states that
cause the blow-up nondeterministic in hybrid automata for deep packet inspection. These
approaches are heuristic. Quasi-deterministic finite-state automata are fully
deterministic devices whose size is bounded in terms of a separator.

\paragraph{Deterministic devices for transducers.}
Subsequential finite-state transducers represent exactly the functions of bounded
variation that are realized by finite-state transducers \cite{mohri_applications}.
Bimachines represent every regular string function deterministically, using a
left-to-right and a right-to-left deterministic automaton \cite{mihov_schulz}, and can be
constructed space-efficiently \cite{gerdjikov_mihov_schulz}. Quasi-subsequential
finite-state transducers, like subsequential ones, read the input once from left to
right and are restricted to functions of bounded variation, but their size is bounded in
terms of a separator.
